\documentclass[12pt,a4paper]{article}
\usepackage[brazilian]{babel}
\usepackage{amsmath,amssymb,amsthm}
\usepackage{geometry}
\usepackage{enumitem}
\usepackage{hyperref}
\usepackage{xcolor}
\usepackage{graphicx}
\usepackage{fancyhdr}
\usepackage{tcolorbox}
\usepackage{totpages}
\usepackage{titlesec}
\usepackage{tasks}

% Configuração global das alternativas com o pacote tasks
\settasks{
  label=\textbf{\Alph*)},
  label-width=1.5em,
  label-offset=0.3em,
  item-indent=0pt,
  column-sep=1.2em
}

% Margens ajustadas para cabeçalho institucional
\geometry{top=3.5cm, bottom=2.5cm, left=2.0cm, right=2.0cm, headheight=2.6cm}

% Paleta de Cores
\definecolor{primary}{RGB}{10, 45, 90}       % Azul Escuro Institucional
\definecolor{secondary}{RGB}{25, 110, 180}   % Azul Destaque
\definecolor{bglight}{RGB}{245, 247, 250}    % Cinza Claro para Caixas
\definecolor{borderlight}{RGB}{210, 220, 230}% Borda Suave
\definecolor{correctgreen}{RGB}{40, 140, 60}  % Verde para Gabarito

% Configuração das Caixas (tcolorbox)
\tcbuselibrary{skins, breakable}

% Configuração de Cabeçalho e Rodapé
\pagestyle{fancy}
\fancyhf{}

% Cabeçalho com Logo e Disciplina
\fancyhead[L]{%
  \includegraphics[height=2cm]{Imagens/logocaruaru.png}%
}
\fancyhead[R]{%
  \small\sffamily
  \textbf{Instituto Federal de Pernambuco -- Campus Caruaru} \\
  \textcolor{secondary}{Disciplina:} Matemática VI \\
  \textcolor{secondary}{Prof.:} Cleibson José da Silva
}

% Rodapé
\fancyfoot[L]{\small\sffamily\color{gray}Lista 6 -- Equações Algébricas1}
\fancyfoot[R]{\small\sffamily\color{gray}Página \thepage\ de \ref{TotPages}}

% Linhas divisórias
\renewcommand{\headrulewidth}{0.8pt}
\renewcommand{\headrule}{{\color{primary}\hrule width\headwidth height \headrulewidth}}
\renewcommand{\footrulewidth}{0.4pt}

\hypersetup{
    colorlinks=true,
    linkcolor=secondary,
    urlcolor=secondary,
    citecolor=primary
}

\begin{document}

% Cabeçalho de Identificação do Estudante
\begin{tcolorbox}[colback=bglight, colframe=primary, boxrule=1pt, arc=3pt]
  \sffamily\small
  \begin{tabular}{@{}p{0.65\textwidth} p{0.3\textwidth}@{}}
    \textbf{Estudante:} \hrulefill & \textbf{Data:} \hrulefill \\[0.25cm]
    \textbf{Curso/Turma:} \hrulefill & \textbf{Nota/Visto:} \hrulefill
  \end{tabular}
\end{tcolorbox}

% Título Principal da Lista
\begin{center}
  {\LARGE\bfseries\color{primary} Lista 6: Equações Algébricas 1} \\[0.1cm]
  {\small\sffamily\color{gray} Equações algébricas, Teorema Fundamental da Álgebra}
\end{center}

\vspace{0.1cm}

 

\begin{enumerate}

    \item \textbf{(FUVEST-SP)} O número de pontos de interseção dos gráficos das funções reais $f(x) = \dfrac{x^2 + 1}{x^2 + 2}$ e $g(x) = \dfrac{x^2 + 4}{x^2 + 3}$ é
   
    \begin{enumerate}
        \item $0$
        \item $1$
        \item $2$
        \item $3$
        \item $4$
    \end{enumerate}

    \item \textbf{(PUC Minas)} Sendo $p(x) = 2x^3 - 5x^2 + 4x - 1$ e $q(x) = 2x^3 - 7x^2 + 7x - 2$, nota-se que $p(1) = q(1) = 0$. A forma mais simples da fração $\dfrac{p(x)}{q(x)}$ é
    \begin{enumerate}
        \item $\dfrac{x + 1}{x - 2}$
        \item $\dfrac{x - 2}{x + 1}$
        \item $\dfrac{x - 1}{x - 2}$
        \item $\dfrac{x - 1}{x + 2}$
        \item $\dfrac{x + 1}{x + 2}$
    \end{enumerate}

    \vspace{0.3cm}

    \item \textbf{(UCS-RS)} Sabe-se que o polinômio $f(x) = x^4 - 4x^3 + 4x^2 - 9$ é divisível por $g(x) = x^2 - 2x + 3$. Se $q(x)$ é o quociente da divisão de $f(x)$ por $g(x)$, quais são as raízes de $q(x)$?
    \begin{enumerate}
        \item $1$ e $-1$
        \item $3$ e $-3$
        \item $1$ e $-3$
        \item $-1$ e $3$
        \item $-1$ e $-3$
    \end{enumerate}

    \vspace{0.3cm}

    \item \textbf{(PUC-SP)} O número de raízes reais do polinômio $p(x) = (x^2 + 1)(x - 1)(x + 1)$ é
    \begin{enumerate}
        \item $0$
        \item $1$
        \item $2$
        \item $3$
        \item $4$
    \end{enumerate}

    \vspace{0.3cm}

    \item \textbf{(PUC Minas)} A equação de terceiro grau cujas raízes são $1$, $2$ e $3$ é
    \begin{enumerate}
        \item $x^3 - 6x^2 + 11x - 6 = 0$
        \item $x^3 - 4x^2 + 4x - 1 = 0$
        \item $x^3 + x^2 + 3x - 5 = 0$
        \item $x^3 + x^2 + 2x + 3 = 0$
        \item $x^3 + 6x^2 - 11x + 5 = 0$
    \end{enumerate}

    \vspace{0.3cm}

    \item \textbf{(UFRN)} Uma das soluções da equação $x^4 - 8x^2 + 16 = 0$ é
    \begin{enumerate}
        \item $-1$
        \item $-2$
        \item $-3$
        \item $-4$
        \item $-5$
    \end{enumerate}

    \vspace{0.3cm}

    \item \textbf{(Cesgranrio)} Sejam $a$ e $b$, respectivamente, a maior e a menor das raízes de $x^4 - 10x^2 + 9 = 0$. A diferença $a - b$ vale
    \begin{enumerate}
        \item $6$
        \item $5$
        \item $4$
        \item $3$
        \item $2$
    \end{enumerate}

    \vspace{0.3cm}

    \item \textbf{(UFRN)} Seja $P(x) = x^3 + 6x^2 - x - 30$. Se $P(2) = 0$, então o conjunto solução de $P(x) = 0$ é
    \begin{enumerate}
        \item $\{-2, -3, -5\}$
        \item $\{2, -3, -5\}$
        \item $\{2, -2, -2\}$
        \item $\{2, 3, 5\}$
        \item $\{2, 6, 30\}$
    \end{enumerate}
\end{enumerate}

\textbf{Gabarito:} 1) A\quad 2) C\quad 3) D\quad 4) C\quad 5) A\quad 6) B\quad 7) A\quad 8) B\quad 
\end{document}